\section{Certified computation and the celestial-mechanics quintic}\label{sec:applications}

\subsection{At least one real root of every real quintic}

Let
\[
 P(x)=a_5x^5+a_4x^4+\cdots+a_0,
 \qquad a_5\ne0.
\]
a polynomial with rational coefficients.  Put
\[
 M=2+\sum_{j<5}\left|\frac{a_j}{a_5}\right|.
\]
For \(x=M\), the leading term dominates the sum of lower terms, and for
\(x=-M\) it has the opposite sign.  More explicitly,
\[
 \left|\sum_{j<5}a_jx^j\right|
 \le |a_5|(M-2)\sum_{j<5}M^j
 <|a_5|M^5.
\]
after replacing the displayed coarse bound by its elementary finite geometric
estimate.  Thus \(P(-M)P(M)<0\).  Exact rational bisection produces a regular
real root with a primitive-recursive modulus.

This proof is the constructive odd-degree theorem needed for quintics.  It does
not invoke the full fundamental theorem of algebra.  To find all real roots,
one may use either of two finite methods:

\begin{enumerate}
\item compute the quartic derivative, isolate its real roots using the checked
      lower-degree factorization identities, and bisect on the resulting
      monotonicity intervals; or
\item compute a signed subresultant/Sturm chain and use its exact variation
      count on rational endpoints.
\end{enumerate}

Both algorithms are primitive recursive for fixed degree five.  Multiple roots
are detected by the polynomial gcd \(\gcd(P,P')\).

\subsection{Certified evaluation of the Bring radical}

For the real Bring radical, bisection has two useful a posteriori certificates.
If a rational interval \([L,U]\) satisfies
\[
 P_a(L)\le0\le P_a(U).
\]
the root lies in the interval and its midpoint error is at most
\((U-L)/2\).  If a single rational candidate \(q\) is available, the sharper
residual estimate~\eqref{eq:residual-bound} gives
\[
 |q-\BR(a)|\le|q^5+q+a|.
\]
Thus a numerical library may use Newton, Halley, or a hypergeometric series for
speed and then certify the answer by one exact rational residual computation.
The executable artifact illustrates this with 80 exact bisection steps for
\(a=3\); it outputs a rational interval of width \(8/2^{80}\).

Inside \(|a|<R\), Theorem~\ref{thm:bring-radius} supplies a series tail bound.
Outside that disk the real branch remains easy: bisection is global because of
strong monotonicity.  Complex branches can be continued by overlapping
rational disks, each equipped with a lower bound for \(|5x^4+1|\) and a
quantitative inverse radius.  The branch history is a finite path certificate.

\subsection{Clearing the Lagrange-point denominators}

Let \(M_S,M_E>0\) be the two primary masses, let \(R>0\) be their separation,
and let \(r>0\) denote the distance of a collinear test mass from the second
primary.  Introduce \(\varepsilon\in\{+1,-1\}\) to encode the two sign
conventions.  The force-balance equation printed on the source page is
\begin{equation}\label{eq:lagrange-force}
 \frac{GM_S}{(R+\varepsilon r)^2}
 +\varepsilon\frac{GM_E}{r^2}
 =\omega^2(R+\varepsilon r).
\end{equation}
The satellite mass has already been cancelled.  Under the nonzero-denominator
guards \(r\ne0\) and \(R+\varepsilon r\ne0\), multiplication by
\(r^2(R+\varepsilon r)^2\) gives
\begin{equation}\label{eq:lagrange-cleared}
 GM_Sr^2+\varepsilon GM_E(R+\varepsilon r)^2
 -\omega^2r^2(R+\varepsilon r)^3=0.
\end{equation}

\begin{theorem}[The collinear equation is quintic]
\label{thm:lagrange-quintic}
If \(\omega\ne0\), the left side of~\eqref{eq:lagrange-cleared} is a polynomial
of degree five in \(r\), with leading coefficient \(-\varepsilon\omega^2\).
Equation~\eqref{eq:lagrange-force} is equivalent to this quintic under the two
denominator guards.
\end{theorem}

\begin{proof}
The last term has degree \(2+3=5\); the other terms have degree at most two.
Its leading monomial is
\(-\omega^2r^2(\varepsilon r)^3=-\varepsilon\omega^2r^5\).
Multiplication by the nonzero common denominator is reversible.  All statements
are polynomial identities and guarded field cancellation.
\end{proof}

Using Kepler's relation
\[
 \omega^2=\frac{G(M_S+M_E)}{R^3}.
\]
a convenient integral scaling of~\eqref{eq:lagrange-cleared} is
\begin{equation}\label{eq:lagrange-kepler}
 \begin{split}
 0={}&M_SR^3r^2+\varepsilon M_ER^3(R+\varepsilon r)^2\\
    &\quad-(M_S+M_E)r^2(R+\varepsilon r)^3.
 \end{split}
\end{equation}
Expanding with \(\varepsilon^2=1\) gives
\begin{equation}\label{eq:lagrange-expanded}
 \begin{split}
 0={}&-\varepsilon(M_S+M_E)r^5
      -3R(M_S+M_E)r^4
      -3\varepsilon R^2(M_S+M_E)r^3\\
    &+(\varepsilon-1)M_ER^3r^2
      +2M_ER^4r+\varepsilon M_ER^5.
 \end{split}
\end{equation}
An overall sign or the replacement \(r\mapsto-r\) changes the signs in a
coefficient table, which explains apparent convention differences among
references; the denominator-cleared equation~\eqref{eq:lagrange-cleared} is the
unambiguous certificate.

\subsection{The Hill-radius approximation}

Let
\[
 \mu=\frac{M_E}{M_S+M_E},
 \qquad s=\frac rR.
\]
For the upper-sign equation \(\varepsilon=1\), dividing
\eqref{eq:lagrange-kepler} by \((M_S+M_E)R^5\) yields the exact relation
\begin{equation}\label{eq:L2-dimensionless}
 \mu(1+2s)=s^3(3+3s+s^2).
\end{equation}
For small \(s\), the dominant balance is \(3s^3=\mu\), giving
\begin{equation}\label{eq:hill-approx}
 r\approx R\sqrt[3]{\frac{M_E}{3(M_S+M_E)}}
 \sim R\sqrt[3]{\frac{M_E}{3M_S}}.
\end{equation}
The lower-sign equation has the same leading balance.

This approximation can be certified without asymptotic notation.  Given a
rational upper bound \(s\le\delta<1/2\), the rational factor
\[
 A(s)=\frac{3+3s+s^2}{1+2s}.
\]
differs from \(3\) by an explicitly bounded rational multiple of \(s\).
Equation~\eqref{eq:L2-dimensionless} is \(\mu=A(s)s^3\).  Rational upper and
lower bounds on \(A\), followed by monotonicity of the cube, give an interval
around \((\mu/3)^{1/3}\).  Alternatively one evaluates the exact quintic at two
rational endpoints and invokes bisection.  Thus every requested relative error
in~\eqref{eq:hill-approx} has a primitive-recursive smallness threshold for
\(\mu\).

\subsection{A general certificate-oriented solver}

The preceding results suggest the following exact workflow for a rational
quintic.

\begin{enumerate}
\item Normalize and factor over \(\Q\).  Solve factors of degree at most four
      with the checked factorization identities.
\item For an irreducible quintic, gather modular factorization data and a
      discriminant.  This often identifies the Galois group.  If the group is
      solvable, construct Lagrange resolvents as in
      Theorem~\ref{thm:solvable-quintic-resolvents}.
\item If the group is nonsolvable, compute a principal and then a certified
      Bring--Jerrard transform.  Evaluate the required Bring branch by series,
      bisection, or analytic continuation, and invert the Tschirnhaus map by a
      subresultant certificate.
\item In every case, substitute the final candidates into the original
      polynomial and verify a complete factorization.  Numerical approximation
      chooses branches but never serves as the final proof.
\end{enumerate}

Every loop has a fixed degree bound and every final claim is an exact polynomial
or finite-table certificate.  This is the form of a quintic solver that is most
natural for \ipra{} and for proof-producing software.
